Human gaze is a promising supervision signal for autonomous driving and robot learning, but only if the measured gaze location is spatially reliable in the conditions where it is recorded. We present a systematic experimental study of a calibration-free wearable eye tracker that proceeds from real-world measurement to controlled investigation and back to driving. First, a gaze-guided, human-in-the-loop pipeline, combining voice-annotated fixations, targeted object detection, automatic target assignment and manual validation, turns natural driving recordings into a benchmark of 152 validated scenes from two qualified participants, with precise 2D fixation-target references, camera--LiDAR target distance, GNSS ego motion and scenario descriptions. Median angular errors of 3.96\dg{} and 1.48\dg{} without correction drop to 0.76\dg{} and 0.83\dg{} with a constant offset estimated from 20\% of the driving scenes, outperforming offsets derived from chessboard calibration (0.88\dg{} and 1.40\dg{}). Second, a controlled study with 216 recordings from six participants shows that accuracy is robust to eccentricity, illumination, indoor/outdoor conditions and time within a session, whereas viewing distance shifts gaze, the lower image region loses precision and smooth pursuit triples residual dispersion. Third, a person-specific causal correction, selected on calibration scenes only, improves on the median-offset baseline in every evaluated split configuration (by 2.5--6.7\% and 7.3--11.8\% for the two participants), with the largest gains in dynamic scenes. Calibration-free wearable gaze is thus a promising foundation for gaze supervision when complemented by sparse in-situ reference observations.
